% ----- CHAUPAI 22 -----

\newpage

\subsection*{\deva{द्वाविंशतितम चौपाई} \(Twenty-Second Chaupai\)}
\addcontentsline{toc}{subsection}{\deva{द्वाविंशतितम चौपाई} \(Twenty-Second Chaupai\) — \deva{सब सुख...}}

\begin{minipage}[t]{0.55\textwidth}
\vspace{0pt}

{\large \linespread{1.5}\selectfont
\deva{सब सुख लहै तुम्हारी सरना।}\\
\deva{तुम रच्छक काहू को डर ना॥}
\par}

\vspace{0.5em}

{\large \linespread{1.5}\selectfont
\telu{సబ సుఖ లహై తుమ్హారీ సరనా।}\\
\telu{తుమ రచ్ఛక కాహూ కో డర నా॥}
\par}

\vspace{0.5em}

{\large \linespread{1.3}\selectfont
\textit{saba sukha lahai tumhārī saranā |}\\
\textit{tuma racchaka kāhū ko ḍara nā ||}
\par}
\end{minipage}%
\hfill
\begin{minipage}[t]{0.42\textwidth}
\vspace{0pt}
\begin{tcolorbox}[anchorbox]
\textbf{Creating Psychological Safety:} A good team captain or leader provides an umbrella of protection so their team can operate without fear (\textit{Dara na}). When you take shelter in good habits and associate with strong, supportive friends, daily anxieties naturally fade away.
\vspace{0.5em}\newline By speaking gently and confidently, he provided immediate psychological refuge to someone trapped in deep anxiety.\newline\null\hfill\href{https://www.valmiki.iitk.ac.in/sloka?field_kanda_tid=5&language=dv&field_sarga_value=34}{\textit{Sundara Kanda, 34:38-40}}
\end{tcolorbox}
\end{minipage}

\vspace{0.5em}
\subsubsection*{\textbf{Visual Rhythm Map:}}
\begin{table}[h!]
\centering
\renewcommand{\arraystretch}{1.2}
\setlength{\tabcolsep}{0pt}
\begin{tabularx}{\textwidth}{|*{16}{>{\centering\arraybackslash\hsize=1.0\hsize}X|}}
\hline
\multicolumn{2}{|>{\centering\arraybackslash\hsize=2.0\hsize}X|}{\textbf{\laghu{स}\laghu{ब}} \par (2)} &
\multicolumn{2}{>{\centering\arraybackslash\hsize=2.0\hsize}X|}{\textbf{\laghu{सु}\laghu{ख}} \par (2)} &
\multicolumn{3}{>{\centering\arraybackslash\hsize=3.0\hsize}X|}{\textbf{\laghu{ल}\guru{है}} \par (3)} &
\multicolumn{5}{>{\centering\arraybackslash\hsize=5.0\hsize}X|}{\textbf{\laghu{तु}\guru{म्हा}\guru{री}} \par (5)} &
\multicolumn{4}{>{\centering\arraybackslash\hsize=4.0\hsize}X|}{\textbf{\laghu{स}\laghu{र}\guru{ना}} \par (4)} \\
\hline
\end{tabularx}
\vspace{-1.5em}
\end{table}

\begin{table}[h!]
\centering
\renewcommand{\arraystretch}{1.2}
\setlength{\tabcolsep}{0pt}
\begin{tabularx}{\textwidth}{|*{16}{>{\centering\arraybackslash\hsize=1.0\hsize}X|}}
\hline
\multicolumn{2}{|>{\centering\arraybackslash\hsize=2.0\hsize}X|}{\textbf{\laghu{तु}\laghu{म}} \par (2)} &
\multicolumn{4}{>{\centering\arraybackslash\hsize=4.0\hsize}X|}{\textbf{\guru{र}\laghu{च्छ}\laghu{क}} \par (4)} &
\multicolumn{4}{>{\centering\arraybackslash\hsize=4.0\hsize}X|}{\textbf{\guru{का}\guru{हू}} \par (4)} &
\multicolumn{2}{>{\centering\arraybackslash\hsize=2.0\hsize}X|}{\textbf{\guru{को}} \par (2)} &
\multicolumn{2}{>{\centering\arraybackslash\hsize=2.0\hsize}X|}{\textbf{\laghu{ड}\laghu{र}} \par (2)} &
\multicolumn{2}{>{\centering\arraybackslash\hsize=2.0\hsize}X|}{\textbf{\guru{ना}} \par (2)} \\
\hline
\end{tabularx}
\vspace{-1.5em}
\end{table}

\vspace{0.5em}
\textbf{ \deva{सरल-संस्कृतम्}:}\\
 \deva{यः भवतः शरणम् आगच्छति, सः सर्वाणि सुखानि लभते।}\\
 \deva{यदा भवान् रक्षकः अस्ति, तदा कस्यापि भयं नास्ति॥}\\

Whoever takes your shelter, obtains all comforts and joys.\\
When you are the protector, there is nothing to fear.


\vspace{1em}
\newpage
\textbf{\deva{प्रतिपदार्थः} (Meaning of each word):}
\begin{table}[h!]
\centering
\renewcommand{\arraystretch}{1.2}
\begin{tabularx}{\textwidth}{@{} l l X X @{}}
\toprule
\textbf{अवधी} & \textbf{संस्कृतम्} & \textbf{English} & \textbf{Grammar \& Notes} \\
\midrule
\deva{सब सुख} / \telu{సబ సుఖ} / \textit{saba sukha}  &  \deva{सर्वाणि सुखानि}  &  All happiness  &   \\
\deva{लहै} / \telu{లహై} / \textit{lahai}  &  \deva{लभते}  &  Obtains / gets  & Verb ending in `-ai` \\
\deva{तुम्हारी} / \telu{తుమ్హారీ} / \textit{tumhārī}  &  \deva{भवतः}  &  Your  &   \\
\deva{सरना} / \telu{సరనా} / \textit{saranā}  &  \deva{शरणम्}  &  Shelter / Refuge  & \\ 
\deva{तुम} \gotchaicon / \telu{తుమ} / \textit{tuma}  &  \deva{त्वम्}  &  You  &   \\
\deva{रच्छक} \gotchaicon / \telu{రచ్ఛక} / \textit{racchaka}  &  \deva{रक्षकः}  &  Protector  & \\ 
\deva{काहू को} / \telu{కాహూ కో} / \textit{kāhū ko}  &  \deva{कस्यापि}  &  For anyone  &   \\
\deva{डर ना} / \telu{డర నా} / \textit{ḍara nā}  &  \deva{भयं नास्ति}  &  No fear  &   \\
\bottomrule
\end{tabularx}
\end{table}

\begin{tcolorbox}[gotchabox]
\begin{itemize}
    \item \textbf{The \deva{क्ष} $\rightarrow$ \deva{च्छ} Consonant Shift:} Look closely at \deva{रच्छक} (Racchaka). The mighty Sanskrit conjunct \deva{क्ष} (kṣa) is very hard on the tongue. In Awadhi, it either softens to \deva{ख} (kha) as in \deva{लखन} (Lakhan for Lakshman), or it softens into the palatal \deva{च्छ} (ccha) as we see here (\deva{रक्षक} $\rightarrow$ \deva{रच्छक}).
\end{itemize}
\end{tcolorbox}

\newpage
